\section{Setup}
\label{sec:setup}

\paragraph{Operator.}
An RGB image $x\in[-1,1]^{3\times256\times256}$ is mapped to $z=(x+1)/2$, and each channel is zero-padded to a $384\times384$ canvas by $P$.
The measurement of a ground-truth image $\xs$ is
\begin{equation}
y=\cA(\xs)+\sigma n,\qquad \cA(x)=\bigl|\cF\,P\,z\bigr|,\qquad n\sim\mathcal N(0,I),
\label{eq:operator}
\end{equation}
with $\cF$ the centered, orthonormal 2-D DFT applied per channel and $\sigma=0.05$.
This is the operator of DPS and DAPS~\cite{chung2023dps,zhang2025daps} (``oversampling 2.0'') and of the solvers we study\claim{C-S1}.
We keep $y$ signed.
With $m=3\cdot384^2$, the residual of a reconstruction is $\rho(x)=\|\cA(x)-y\|_2/\sqrt m$.
The support occupies $2/3$ of the canvas per dimension, outside the classical uniqueness condition for generic images~\cite{bruck1979ambiguity,barnett2020geometry}.
We assume that $\xs$ is determined by $y$ up to the symmetries below and to images that fit the data within the noise\claim{C-S3}.

\paragraph{The orbit, and what selection cannot see.}
For $b\in\{0,1\}^3$, let $R_b$ rotate color channel $c$ by $180^\circ$ inside the $256\times256$ frame whenever $b_c=1$.
On the even canvas this is index negation followed by a one-pixel cyclic shift, so $\cA(R_bx)=\cA(x)$.
This is the classical twin ambiguity~\cite{hayes1982reconstruction,fienup1986stagnation}, per channel; Zhang~\etal display it for DAPS~\cite[App.~F]{zhang2025daps}\claim{C-T1}.
We call the eight images $R_bx$ the \emph{orbit} of $x$; $R_{111}$ rotates the whole image.
Two consequences matter here.
First, any rule that selects among runs using only $\cA(\hat x_1),\dots,\cA(\hat x_K)$ and $y$, including the residual, returns the same index when a run is replaced by any member of its orbit; so do likelihood-only particle weights\claim{C-T2}.
Second, the members of an orbit are pixel permutations with the same likelihood, so $p(R_bx\mid y)/p(x\mid y)=\pi(R_bx)/\pi(x)$ for a prior density $\pi$.
Within an orbit, the posterior ranks images by the prior alone, and any orientation test that orders orbit members like $\pi$ is the Bayes choice.
The supplement bounds why samplers prefer upright images\claim{C-T4}.

\paragraph{Run labels.}
PSNR is computed on 8-bit RGB against $\xs$, with a threshold $\theta$ of $25$\,dB on FFHQ and $20$\,dB on ImageNet.
A run $\hat x$ is
\good if $\mathrm{PSNR}(\hat x)\ge\theta$;
an \orbit error if it is not \good but $\max_b\mathrm{PSNR}(R_b\hat x)\ge\theta$ (\emph{global} if the best member is $R_{111}$, \emph{partial} otherwise);
\cons if it is neither but $\rho(\hat x)\le\tau$, \ie it fits the measurement as well as the solver's \good runs do;
and \incons otherwise.
Here $\tau$ is the $95$th percentile of $\rho$ over the \good runs of the same solver and cohort, fixed before the analyses reported here\claim{C-S4}.
Labels are used only to score policies.
The comparisons do not hinge on $\theta$: from $3$\,dB below to $2$\,dB above it, $\Eone$ stays within $2.8$ points of the orbit oracle for SITCOM, and choosing among PC-DAPS particles keeps its gain (supplement)\claim{C-V1}.

\paragraph{Policies.}
Given $K$ runs of one image, we compare
the \emph{first run};
the residual pick $\Eo=\argmin_k\rho(\hat x_k)$;
$\Eone$, which applies an orientation test to the $\Eo$ pick and returns it or its rotation $R_{111}$;
the ground-truth oracle $\Oraw=\argmax_k\mathrm{PSNR}(\hat x_k)$;
the orbit oracle $\Oorb$, which also maximizes over the eight members of each orbit;
and the \emph{per-metric oracle}, which takes the best PSNR, SSIM and LPIPS independently, as the reference code does\claim{C-P1}.
Only the first run, $\Eo$, $\Eone$ and the protocol of \cref{sec:protocol} are available at test time.
For a policy and a set of images, \good denotes the share of images whose output is \good.

\paragraph{Data and solvers.}
\emph{FFHQ}~\cite{karras2019stylegan}: $180$ images in two disjoint cohorts, with the $256\times256$ denoiser of DPS~\cite{chung2023dps}.
Cohort~A has $100$ images drawn at random from the validation split.
Cohort~B has $80$ images from an earlier screen enriched for hard cases; $71$ of them come from the training split, which the denoiser may have seen, so we report the cohorts separately; the conclusions hold on cohort~A alone\claim{C-S7}.
\emph{ImageNet}~\cite{deng2009imagenet}: the first $100$ images of the test archive distributed with DAPS, the default evaluation set of DAPS, FAST-DIPS and PC-DAPS, with the unconditional denoiser of Dhariwal and Nichol~\cite{dhariwal2021diffusion}, trained on the training split; all our ImageNet images are validation images.
\emph{ImageNet$'$}: $100$ further ImageNet validation images that no earlier experiment touched (\cref{sec:protocol}).
Each image has one measurement, shared by all solvers and runs.
Each run uses a fresh seed, with no warm start and no tuning on other runs, so the runs of an image are exchangeable; we found no drift with the run index and no dependence between consecutive runs\claim{C-F4}\claim{C-V3}.
\emph{Solvers}: DAPS~\cite{zhang2025daps} in two configurations, named by their denoiser calls per run: DAPS-4K, the configuration of its published phase-retrieval results ($400$ annealing steps with $10$ ODE steps each), and, on FFHQ, DAPS-2K ($400$ steps with $5$ ODE steps each)\claim{C-S8};
SITCOM~\cite{alkhouri2025sitcom}, through its authors' phase-retrieval code with the published settings;
FAST-DIPS~\cite{kim2026fastdips}, with its published command;
and PC-DAPS~\cite{kocabas2026pcdaps}, with its published phase-retrieval configuration (\cref{sec:audit}).
DAPS and SITCOM weight the data term $50$ and $25$ times more than the Gaussian likelihood at $\sigma=0.05$, and SITCOM adds a gradient step on the residual after each ODE step.
FAST-DIPS returns a denoised estimate after a hard measurement projection\claim{C-S5}.
A \emph{panel} is one solver on one image set, with a fixed number of runs or budget per image.
